% SOURCE_AUTHORITY: sga5_fr_workpass.tex, lines 7160--7177 (LF-normalized unit).
% SOURCE_SHA256: 791F4EFFC5E02832D5D77ED03518C8156D6F07E4C8238B03545DB93D883FBB28
\begin{prooftext}[Demonstração]
\begin{enumerate}[label=(\alph*)]
\item As condições do enunciado são suficientes: se $X$ verifica (MLAR), é AR-isomorfo ao seu sistema projetivo das imagens universais $X'$. Por outro lado, a hipótese feita sobre $X'$ significa que o sistema projetivo
\[
\bigl(X'_{n+r}/J^{n+1}X'_{n+r}\bigr)_{n\ge0},
\]
com os morfismos de transição evidentes, é $J$-ádico. Ora, é claro que este sistema projetivo é AR-isomorfo (2.4.6) a $X'$.
\item As condições do enunciado são necessárias: se $X$ é AR-$J$-ádico, é AR-isomorfo a um sistema projetivo estrito e, portanto, verifica (MLAR). O sistema projetivo $X'$ das imagens universais de $X$ é AR-isomorfo a $X$ e, portanto, AR-$J$-ádico. Para demonstrar o segundo ponto, reduzimo-nos, pois, a verificar que, se $X$ é um sistema projetivo AR-$J$-ádico e estrito, existe um inteiro $r$ tal que o sistema projetivo $\ell_r(X)$ definido por
\[
\ell_r(X)_n=X_{n+r}/J^{n+1}X_{n+r}
\]
seja $J$-ádico. Para isso, seja $u:Y\to X$ um AR-isomorfismo, sendo $Y$ $J$-ádico. Como o núcleo $T$ de $u$ é AR-nulo, seja $r$ um inteiro tal que o morfismo canônico $T[r]\to T$ seja nulo. É imediato que o morfismo canônico
\[
\ell_r(u):\ell_r(Y)\longrightarrow\ell_r(X)
\]
é um isomorfismo, de onde segue-se a afirmação.
\end{enumerate}
\end{prooftext}
